\section{Introduction}

\emph{``What is the average height for a male?''} --- given a formal and a casual answer with verified identical content (NLI entailment $= 0.993$), GPT-4o rates the casual variant \emph{higher} on a 1--10 Likert scale (8 vs.\ 5), yet consistently selects the formal variant in pairwise comparison.
Same judge, same texts, opposite conclusions.
We term this the \emph{Likert Blindspot}---an operational descriptor for the measurement-mode divergence where pairwise comparison detects statistically significant style sensitivity that Likert scoring fails to detect on the same data, without committing to a specific causal mechanism.
This divergence admits multiple non-exclusive explanations---Likert measurement attenuation, pairwise format amplification \citep{jeong2025comparative}, or construct-level differences between absolute and comparative evaluation---which our controlled design bounds but cannot definitively adjudicate (\S\ref{sec:likert_pairwise}).
Under construct divergence, pairwise and Likert modes may measure distinct constructs---comparative sensitivity versus absolute quality judgment---in which case the divergence reflects measurement scope rather than measurement failure, with different practical implications for audit design (\S\ref{sec:likert_pairwise}).
This example reflects a consistent pattern: surface-level style features alone explain most of the variance in judge scores \citep{murugadoss2025sos}, with 76--92\% of preferences originating from style cues \citep{li2026preference,soumik2026judging}.

LLM judges are central infrastructure for NLP evaluation and alignment training (RLHF, DPO), where judge preferences serve as reward signals; any systematic bias propagates through both evaluation leaderboards and model development pipelines.
Despite this, existing approaches address individual requirements (multi-dimensional style analysis \citep{casse2026concise,wang2026fairjudge}, interventional design \citep{reber2025rate}, population-level correction \citep{dubois2025lc}), but no single framework integrates per-axis decomposition, interventional identification, per-instance diagnostics, and black-box applicability.
The classical psychometric literature establishes that pairwise comparison yields higher discriminability than absolute rating \citep{thurstone1927law,david1988method}, and our power analysis (Figure~\ref{fig:power_curve}) is quantitatively consistent: Likert requires $4$--$206\times$ the sample size.
However, the LLM judge setting introduces content-specific structure that complements the classical framework: (a)~the Blindspot exhibits \emph{per-axis structure}---verbosity is style-dominated (Style OR $8$--$13$) while formality is quality-confounded---complementing Thurstone's format-general discriminability gradient (Table~\ref{tab:gee}); (b)~GEE decomposition separates the \emph{sources} of discriminability into style-specific and quality-correlated components, not standard practice in LLM judge evaluation; (c)~the divergence manifests as a near-complete binary split (all seven measurement methods detect via pairwise, all fail via Likert; \S\ref{sec:likert_pairwise}), consistent with classical predictions when accounting for effect-size-dependent power (Figure~\ref{fig:power_curve}).

We propose Controlled Style Decomposition (CSD), combining NLI-gated content control, per-axis style decomposition, and position-controlled pairwise comparison.
Inspired by interventional text analysis \citep{veitch2020adapting}, CSD operates in three stages: style-controlled rewriting along a single axis, NLI entailment verification, and bias measurement via forced-choice comparison.
CSD requires no access to judge internals and applies post-hoc to any black-box judge.
Our experiments demonstrate the \emph{Likert Blindspot} across four target judges from independent providers (GPT-4o, Qwen3-32B, Llama-3.3-70B, Gemini-2.5-Flash), a fifth judge (Claude Sonnet~4) with benchmark-dependent sensitivity, three benchmarks, and two primary style axes, with exploratory register evidence.
On formality data, Likert yields non-significant ATE ($p = 0.332$, $0.908$), while BT detects 63.0\% (GPT-4o, $p = 0.0004$) and 66.5\% (Qwen3-32B, $p < 0.0001$) formal preferences; verbosity shows 77.3\% BT ($p < 0.0001$) invisible to Likert ($p = 0.179$; Figure~\ref{fig:forest_plot}).

Our contributions are both methodological and empirical:
\begin{itemize}
    \item \textbf{CSD}, a black-box diagnostic integrating NLI-gated content control, per-axis decomposition, per-instance diagnostics, and black-box applicability---capabilities addressed individually by prior work but not jointly (\S\ref{sec:method}).
    \item \textbf{Systematic cross-format comparison} across four judges, three benchmarks, and two axes with position control---the first at this scale---revealing measurement-mode divergence as a consistent pattern rather than an isolated anomaly: complete at moderate bias ($d \to 0$ at $n = 355$), ${>}60\%$ signal loss at stronger effects (Table~\ref{tab:main_results}).
    \item \textbf{GEE conditional decomposition} separating style-specific from quality-correlated components---distinct from factor analysis or SEM in targeting observed binary predictors rather than latent variables---revealing qualitatively distinct confound structures per axis (Table~\ref{tab:gee}).
\end{itemize}

\begin{figure}[t]
\centering
\includegraphics[width=\columnwidth]{figures/pipeline_diagram.pdf}
\caption{CSD pipeline overview. Numbers indicate sample attrition at each stage.}
\label{fig:pipeline}
\end{figure}
